
\documentclass[11pt,letterpaper,oneside]{report}
\usepackage[T1]{fontenc}
\usepackage[utf8]{inputenc}
\usepackage{lmodern}
\usepackage{inconsolata}
\usepackage{geometry}
\geometry{margin=0.89in,headheight=16pt,headsep=19pt,footskip=26pt}
\usepackage{xcolor}
\definecolor{navy}{HTML}{173758}
\definecolor{teal}{HTML}{157B89}
\definecolor{pale}{HTML}{E9F0F3}
\definecolor{codegray}{HTML}{F5F7F9}
\definecolor{muted}{HTML}{667785}
\usepackage{booktabs,longtable,tabularx,array,calc}
\usepackage{amsmath,amssymb}
\usepackage{graphicx}
\usepackage{microtype}
\usepackage{enumitem}
\usepackage{fancyhdr}
\usepackage{titlesec}
\usepackage{listings}
\usepackage{upquote}
\usepackage{xurl}
\usepackage{hyperref}
\hypersetup{colorlinks=true,linkcolor=navy,urlcolor=teal,bookmarksopen=true,pdfauthor={Software Build Engineering Documentation Suite}}
\urlstyle{same}
\setcounter{tocdepth}{2}
\setcounter{secnumdepth}{2}
\setlength{\parindent}{0pt}
\setlength{\parskip}{5pt plus 1pt}
\setlength{\emergencystretch}{3em}
\setlist[itemize]{leftmargin=*,topsep=4pt,itemsep=2pt}
\setlist[enumerate]{leftmargin=*,topsep=4pt,itemsep=2pt}
\titleformat{\chapter}[display]{\normalfont\bfseries\color{navy}}{\fontsize{11}{12}\selectfont CHAPTER \thechapter}{7pt}{\fontsize{23}{28}\selectfont}
\titlespacing*{\chapter}{0pt}{-6pt}{19pt}
\titleformat{\section}{\Large\bfseries\color{navy}}{\thesection}{0.65em}{}
\titleformat{\subsection}{\large\bfseries\color{teal}}{\thesubsection}{0.65em}{}
\pagestyle{fancy}
\fancyhf{}
\fancyhead[L]{\small\textcolor{muted}{PROFESSIONAL CMAKE ENGINEERING}}
\fancyhead[R]{\small\textcolor{muted}{\nouppercase{\leftmark}}}
\fancyfoot[C]{\small\textcolor{muted}{\thepage}}
\renewcommand{\headrulewidth}{0.3pt}
\fancypagestyle{plain}{\fancyhf{}\fancyfoot[C]{\small\thepage}\renewcommand{\headrulewidth}{0pt}}
\lstdefinelanguage{cmake}{sensitive=false,morekeywords={cmake_minimum_required,project,add_library,add_executable,target_compile_features,target_include_directories,target_link_libraries,target_sources,target_compile_options,target_compile_definitions,target_link_options,add_subdirectory,include,if,else,elseif,endif,foreach,endforeach,function,endfunction,macro,endmacro,set,option,find_package,find_program,add_custom_command,add_custom_target,install,configure_file,configure_package_config_file,write_basic_package_version_file,set_target_properties,set_tests_properties,add_test,message,FetchContent_Declare,FetchContent_MakeAvailable,return,cmake_parse_arguments},morecomment=[l]{\#},morestring=[b]"}
\lstdefinelanguage{sh}{morekeywords={cmake,ctest,cpack,find,sha256sum,clang-tidy,echo,cp,make,ninja,git},morecomment=[l]{\#},morestring=[b]"}
\lstdefinelanguage{yaml}{morecomment=[l]{\#},morestring=[b]",morekeywords={name,on,permissions,jobs,steps,uses,run,with,env,matrix,include,if}}
\lstdefinelanguage{text}{}
\lstset{basicstyle=\ttfamily\footnotesize,keywordstyle=\color{navy}\bfseries,commentstyle=\color{muted},stringstyle=\color{teal},backgroundcolor=\color{codegray},frame=single,rulecolor=\color{pale},breaklines=true,breakatwhitespace=false,showstringspaces=false,columns=fullflexible,keepspaces=true,tabsize=2,xleftmargin=5pt,xrightmargin=5pt,aboveskip=9pt,belowskip=9pt}
\providecommand{\tightlist}{\setlength{\itemsep}{0pt}\setlength{\parskip}{0pt}}
\providecommand{\passthrough}[1]{#1}
\providecommand{\real}[1]{#1}

\begin{document}
\hypersetup{pageanchor=false}\begin{titlepage}
\color{navy}
\vspace*{1.15in}
{\Large\bfseries PROFESSIONAL CMAKE\\[7pt]SOFTWARE BUILD ENGINEERING\par}
\vspace{0.32in}\noindent\textcolor{teal}{\rule{0.9\linewidth}{2.3pt}}\par
\vspace{0.45in}
{\parbox{0.95\textwidth}{\raggedright\fontsize{24}{29}\selectfont\bfseries Prerequisite-Based CMake Curriculum\par}}
\vspace{0.3in}
{\Large\color{teal}Twelve weeks, labs and capstone\par}
\vfill
{\large DOCUMENT 08 OF 08\par}
\vspace{0.1in}
{\small Companion synthesis to Craig Scott, \emph{Professional CMake: A Practical Guide}, 20th Edition (2025).\\
Original learning/engineering material. Not affiliated with or endorsed by the author.\par}
\vspace{0.2in}
{\small 17 October 2026\par}
\end{titlepage}
\hypersetup{pageanchor=true}\clearpage\pagenumbering{roman}\tableofcontents\clearpage\pagenumbering{arabic}

\textbf{Training objective:} develop CMake engineers who can design,
debug, secure, package, and maintain portable, testable C/C++ build
systems. \textbf{Duration:} 12 weeks, approximately 8--10 hours/week,
96--120 guided and independent hours. \textbf{Reference:}
\emph{Professional CMake}, 20th Edition, 44 chapters in seven parts.
\textbf{{[}BOOK-GROUNDED{]}} the original chapter sequence supplies the
technical coverage. \textbf{{[}ENGINEERING EXTENSION{]}} pacing,
prerequisite graph, skills assessment, capstone and enterprise risk
controls are original training design decisions.

\chapter{Entry requirements and learner
tracks}\label{entry-requirements-and-learner-tracks}

\textbf{Required:} working C/C++ compilation and linking knowledge,
command-line shell usage, Git basics and familiarity with at least one
compiler. \textbf{Recommended:} library linking, static versus shared
binaries, unit testing, debugging symbol use, basic Linux/Windows
filesystem semantics and JSON/YAML. There is no requirement to know
advanced CMake before entering this path.

\begin{longtable}[]{@{}
  >{\raggedright\arraybackslash}p{(\columnwidth - 6\tabcolsep) * \real{0.2500}}
  >{\raggedright\arraybackslash}p{(\columnwidth - 6\tabcolsep) * \real{0.2500}}
  >{\raggedright\arraybackslash}p{(\columnwidth - 6\tabcolsep) * \real{0.2500}}
  >{\raggedright\arraybackslash}p{(\columnwidth - 6\tabcolsep) * \real{0.2500}}@{}}
\toprule\noalign{}
\begin{minipage}[b]{\linewidth}\raggedright
Path
\end{minipage} & \begin{minipage}[b]{\linewidth}\raggedright
Entry knowledge
\end{minipage} & \begin{minipage}[b]{\linewidth}\raggedright
Emphasis
\end{minipage} & \begin{minipage}[b]{\linewidth}\raggedright
Specialization outcome
\end{minipage} \\
\midrule\noalign{}
\endhead
\bottomrule\noalign{}
\endlastfoot
Core developer & C++ build basics & target APIs, local builds, tests &
maintain a modern target graph \\
Build platform engineer & multi-platform tools & toolchains, presets,
packaging & own a CI build matrix \\
Embedded developer & cross toolchains & sysroots, ABIs, device
verification & deploy ARM/Android artifacts \\
DevSecOps/AppSec & CI and threat modeling & code scanning, provenance,
dependency governance & validate secure release controls \\
\end{longtable}

\chapter{Dependency graph}\label{dependency-graph}

\begin{lstlisting}
C/C++ compiler & linker basics
             |
        [W1 CMake lifecycle]
             |
        [W2 language/scopes] --> [W3 targets & properties]
                                    |
                         +----------+---------+
                         |                    |
                   [W4 advanced]        [W6 testing]
                         |                    |
                  [W5 toolchains]      [W7 analysis]
                         +----------+---------+
                                    |
                              [W8 deps]
                                    |
                              [W9 install]
                                    |
                            [W10 organization]
                                    |
                          [W11 DevSecOps]
                                    |
                           [W12 capstone]
\end{lstlisting}

Modules can be taken in parallel only after their explicit
target-contract prerequisites are satisfied. The graph is a teaching
design; it does not imply the book prescribes 12 weeks.

\chapter{Week-by-week learning plan}\label{week-by-week-learning-plan}

\section{Week 1. CMake lifecycle and minimal
targets}\label{week-1.-cmake-lifecycle-and-minimal-targets}

\textbf{Book mapping:} Part I, Chapters 1--5. \textbf{Prerequisites:}
compiler and linker command lines. \textbf{Objectives:} differentiate
configure/generate/build/test/package, use out-of-source builds, declare
executable and library targets, and run a basic CTest. \textbf{Lab:}
implement a two-target CLI, two configurations and a staged install.
\textbf{Acceptance:} \passthrough{\lstinline!cmake -S!},
\passthrough{\lstinline!cmake --build!}, \passthrough{\lstinline!ctest!}
and \passthrough{\lstinline!cmake --install!} succeed from a clean
checkout; no generated files are tracked in source.

\section{Week 2. Language, variables, structure and helper
functions}\label{week-2.-language-variables-structure-and-helper-functions}

\textbf{Book mapping:} Chapters 6--9 and 12--13. \textbf{Prerequisites:}
Week 1. \textbf{Objectives:} distinguish cache versus local variables,
functions versus macros, policy baselines, module inclusion and
subdirectory scopes. \textbf{Lab:} create a parametrized helper for
registering test executables using
\passthrough{\lstinline!cmake\_parse\_arguments!}; prove meaningful
error handling. \textbf{Acceptance:} demonstrates normal/cache shadowing
and correct behavior with a space-containing source path.

\section{Week 3. Target properties and
interfaces}\label{week-3.-target-properties-and-interfaces}

\textbf{Book mapping:} Chapters 10--11, 14--16. \textbf{Prerequisites:}
Week 2. \textbf{Objectives:} express compile/link requirements at target
granularity, use generator expressions and inspect build commands.
\textbf{Lab:} design three targets with one PUBLIC and one PRIVATE
dependency. \textbf{Acceptance:} external consumer compiles with no
manually supplied include paths; Debug-only definition is absent from
Release.

\section{Week 4. Advanced build graphs and library
design}\label{week-4.-advanced-build-graphs-and-library-design}

\textbf{Book mapping:} Chapters 17--23. \textbf{Prerequisites:} Week 3.
\textbf{Objectives:} differentiate
STATIC/SHARED/INTERFACE/OBJECT/IMPORTED targets, generation-time custom
tasks, file APIs, library ABI and language requirements. \textbf{Lab:}
generate a version header, build a shared library and test symbol-export
choices. \textbf{Acceptance:} no spurious rebuild on a no-op invocation;
documented shared-library version strategy.

\section{Week 5. Cross-compilation, platforms and
performance}\label{week-5.-cross-compilation-platforms-and-performance}

\textbf{Book mapping:} Chapters 24--26. \textbf{Prerequisites:} Week 4.
\textbf{Objectives:} design a toolchain file, handle sysroot search,
explain try\_run limitations and benchmark compilation options.
\textbf{Lab:} make a cross-compilation plan (actual hardware optional),
and benchmark precompiled headers or compiler cache.
\textbf{Acceptance:} platform support matrix and five-sample benchmark
including build/test regression check.

\section{Week 6. CTest fundamentals and resource
management}\label{week-6.-ctest-fundamentals-and-resource-management}

\textbf{Book mapping:} Chapters 27--30. \textbf{Prerequisites:} Week 3.
\textbf{Objectives:} register tests, use fixtures/resources/timeouts,
integrate a test framework and understand separate build-and-test
projects. \textbf{Lab:} add unit and integration suites, resource lock
and JUnit report. \textbf{Acceptance:} expected failure is captured;
parallel runs are deterministic.

\section{Week 7. Static, dynamic and runtime
analysis}\label{week-7.-static-dynamic-and-runtime-analysis}

\textbf{Book mapping:} Chapters 31--33. \textbf{Prerequisites:} Weeks
5--6. \textbf{Objectives:} use compiler databases, clang-tidy, sanitizer
builds, coverage and dashboard reporting. \textbf{Lab:} isolated
ASan/UBSan build with an intentional negative canary; add a static
check. \textbf{Acceptance:} report demonstrates that a known invalid
condition is detected and that the normal build remains unaffected.

\section{Week 8. Dependency resolution and
control}\label{week-8.-dependency-resolution-and-control}

\textbf{Book mapping:} Chapters 34, 38--39, 41. \textbf{Prerequisites:}
Weeks 3 and 5. \textbf{Objectives:} distinguish find\_package,
FetchContent, ExternalProject and dependency providers. \textbf{Lab:}
implement one controlled external library discovery and contrast with a
pinned source dependency in an architecture record. \textbf{Acceptance:}
release scenario works with approved local dependencies; acquisition
policy includes source, integrity and license.

\section{Week 9. Installation, packaging and SDK
consumption}\label{week-9.-installation-packaging-and-sdk-consumption}

\textbf{Book mapping:} Chapters 35--37 and 40. \textbf{Prerequisites:}
Week 8. \textbf{Objectives:} create GNUInstallDirs installation rules,
exported targets, relocatable config packages, archive and platform
package plan. \textbf{Lab:} install reference library to a custom prefix
and build a completely separate consumer. \textbf{Acceptance:} consumer
still works after moving the installed prefix; package content manifest
matches expected public files.

\section{Week 10. Presets and project
architecture}\label{week-10.-presets-and-project-architecture}

\textbf{Book mapping:} Chapters 42--43, optional 44 for Qt.
\textbf{Prerequisites:} Weeks 4 and 9. \textbf{Objectives:} structure a
larger project, distinguish normal build and superbuild, implement
developer/CI presets and evaluate GUI framework deployment.
\textbf{Lab:} refactor multiple targets into owned subdirectories and
issue two architecture decisions. \textbf{Acceptance:} reusable library
remains safe as a subproject; two presets and clean consumer succeed.

\section{Week 11. Secure CI/CD and
reproducibility}\label{week-11.-secure-cicd-and-reproducibility}

\textbf{Book mapping:} synthesis of Chapters 24, 26--33 and 35--43.
\textbf{Prerequisites:} Weeks 7--10. \textbf{{[}ENGINEERING EXTENSION{]}
Objectives:} implement GitHub Actions or equivalent CI, analysis
coverage, dependency inventory, CodeQL evaluation, artifact hashes and
release policy. \textbf{Lab:} define platform matrix, PR gates and a
release staging workflow. \textbf{Acceptance:} preserved evidence for
pass/fail status, tests, inventories, platform gaps and trust
boundaries.

\section{Week 12. Capstone design, defense and
remediation}\label{week-12.-capstone-design-defense-and-remediation}

\textbf{Book mapping:} all seven parts. \textbf{Prerequisites:} Weeks
1--11. \textbf{Objectives:} demonstrate an independently consumable
software product and defend key technical decisions. \textbf{Lab:}
submit source, build presets, package config, test suite, analysis
records, ADRs and CI design. \textbf{Acceptance:} scoring rubric below;
independent reviewer reproduces from a clean source checkout.

\chapter{Weekly review sequence}\label{weekly-review-sequence}

Each week includes: 90 minutes of guided reading, 60 minutes of concept
recall, 3 hours of hands-on lab work, 90 minutes of
troubleshooting/refactoring, 60 minutes of test and artifact review, and
optional additional hours for platform specialization. Learners maintain
a build log containing revision, compiler/toolchain versions, exact
commands and expected versus actual results.

\chapter{Capstone specification: Portable library and CLI
SDK}\label{capstone-specification-portable-library-and-cli-sdk}

\textbf{System:} build \passthrough{\lstinline!bsuite\_core!} (library),
\passthrough{\lstinline!bsuite\_cli!} (executable), unit and integration
tests, an independent consumer, a package-config export and staged CPack
archive. The deliverable should include at least three distinct build
configurations or platform scenarios and one instrumented analysis path.

\section{Functional requirements}\label{functional-requirements}

\begin{longtable}[]{@{}
  >{\raggedright\arraybackslash}p{(\columnwidth - 4\tabcolsep) * \real{0.3333}}
  >{\raggedright\arraybackslash}p{(\columnwidth - 4\tabcolsep) * \real{0.3333}}
  >{\raggedright\arraybackslash}p{(\columnwidth - 4\tabcolsep) * \real{0.3333}}@{}}
\toprule\noalign{}
\begin{minipage}[b]{\linewidth}\raggedright
ID
\end{minipage} & \begin{minipage}[b]{\linewidth}\raggedright
Requirement
\end{minipage} & \begin{minipage}[b]{\linewidth}\raggedright
Verification
\end{minipage} \\
\midrule\noalign{}
\endhead
\bottomrule\noalign{}
\endlastfoot
FR-01 & configure from a clean source tree with documented CMake minimum
& clean configure transcript \\
FR-02 & build library and CLI via target dependencies & build graph and
\passthrough{\lstinline!--verbose!} evidence \\
FR-03 & propagate public API usage requirements & in-tree and installed
consumer build \\
FR-04 & execute deterministic tests with CTest & JUnit and exit
status \\
FR-05 & install headers, binaries and versioned CMake package & stage
manifest \\
FR-06 & relocate installation without rebuilding producer & new-prefix
consumer smoke \\
FR-07 & demonstrate isolated static/dynamic analysis path & tool report
/ canary \\
FR-08 & compare at least one performance intervention & reproducible
benchmark sheet \\
FR-09 & produce a release archive and inventory & package contents and
SHA-256 \\
FR-10 & justify dependency source and CI permissions & architecture
decisions / threat model \\
\end{longtable}

\section{Nonfunctional requirements}\label{nonfunctional-requirements}

The build must not depend on untracked local paths; no secrets may be
committed; supported operating systems and compilers must be declared;
the default project must not run network downloads. Any third-party
library or optional Qt feature is controlled by a separate documented
option and is not required for the base capstone.

\chapter{Assessment rubric (100
points)}\label{assessment-rubric-100-points}

\begin{longtable}[]{@{}
  >{\raggedright\arraybackslash}p{(\columnwidth - 4\tabcolsep) * \real{0.3333}}
  >{\raggedleft\arraybackslash}p{(\columnwidth - 4\tabcolsep) * \real{0.3333}}
  >{\raggedright\arraybackslash}p{(\columnwidth - 4\tabcolsep) * \real{0.3333}}@{}}
\toprule\noalign{}
\begin{minipage}[b]{\linewidth}\raggedright
Category
\end{minipage} & \begin{minipage}[b]{\linewidth}\raggedleft
Points
\end{minipage} & \begin{minipage}[b]{\linewidth}\raggedright
Passing evidence
\end{minipage} \\
\midrule\noalign{}
\endhead
\bottomrule\noalign{}
\endlastfoot
CMake target architecture & 20 & correct PUBLIC/PRIVATE/INTERFACE
requirements \\
Portable configuration and presets & 15 & clean configure, Debug/Release
separation \\
Library installation and consumption & 20 & exported package and
relocated consumer \\
Test strategy and execution & 15 & CTest/JUnit, negative and isolation
checks \\
Security and dependency governance & 15 & traceable inventory,
appropriate trust boundaries \\
Performance and diagnostics & 10 & repeatable measurement and fault
investigation \\
Documentation and defense & 5 & clear ADRs, commands, risks and scope \\
\end{longtable}

\textbf{Passing:} at least 80/100 overall, with 12/20 in architecture
and 12/20 in installation/consumption; no critical secrets or
release-policy violations. \textbf{Distinction:} 90+/100 with repeatable
multi-platform evidence and an approved dependency acquisition process.

\chapter{Skill checkpoints and
remediation}\label{skill-checkpoints-and-remediation}

\begin{longtable}[]{@{}
  >{\raggedright\arraybackslash}p{(\columnwidth - 4\tabcolsep) * \real{0.3333}}
  >{\raggedright\arraybackslash}p{(\columnwidth - 4\tabcolsep) * \real{0.3333}}
  >{\raggedright\arraybackslash}p{(\columnwidth - 4\tabcolsep) * \real{0.3333}}@{}}
\toprule\noalign{}
\begin{minipage}[b]{\linewidth}\raggedright
Checkpoint
\end{minipage} & \begin{minipage}[b]{\linewidth}\raggedright
Common mistake
\end{minipage} & \begin{minipage}[b]{\linewidth}\raggedright
Targeted remediation
\end{minipage} \\
\midrule\noalign{}
\endhead
\bottomrule\noalign{}
\endlastfoot
W3 & global include path leak & remove global directory and correct
target PUBLIC interface \\
W5 & host dependency during cross build & specify sysroot/root-path mode
and inspect discovered package \\
W7 & sanitizer not on final link & instrument executable and dependency
closure coherently \\
W9 & build-tree path in installed target & fix export usage requirements
and rebuild independent consumer \\
W11 & CI passes while scanner skipped & enforce expected status and
explicit failure on unavailable tool \\
\end{longtable}

\chapter{Instructor and reviewer
guide}\label{instructor-and-reviewer-guide}

Review artifacts at their \textbf{reproducible execution boundary}, not
solely their source appearance. Run the exact commands in a clean
workspace; detect hidden environmental state, compare installed versus
in-tree interfaces, and demand a failing regression before confirming a
purported defect fix. Clearly separate tested behavior from unverified
operating systems, packaging generators or security controls.

\section{Suggested extension
projects}\label{suggested-extension-projects}

\textbf{Embedded:} ARM Linux CMake toolchain and emulator/target smoke
tests. \textbf{Qt:} Qt6 Widgets app with bundle, plugin and deployment
review. \textbf{Large-scale SDK:} multi-library export and semantic
version contract. \textbf{AppSec:} separate CodeQL C/C++ workflow with
explicit mode decision, issue triage and SBOM generation external to
CMake.

\chapter{Reading and reference index}\label{reading-and-reference-index}

Textbook Parts I--VII provide the chapter-based sequence.
Cross-reference this suite's \textbf{Technical Reference} for commands,
\textbf{Implementation Handbook} for executable artifacts, \textbf{Build
Architecture} for ADRs, \textbf{Testing and Quality} for evidence gates,
\textbf{CI/CD and DevSecOps} for software supply-chain controls and
\textbf{Cross-Platform and Performance} for tooling/benchmark methods.
Official docs: \url{https://cmake.org/cmake/help/latest/} .

\end{document}
